\documentclass[aps,prb,twocolumn,superscriptaddress,nobibnotes]{revtex4-2} % APS journal style
\usepackage[a4paper,margin=0.8in]{geometry}
\usepackage{graphicx}
\usepackage{amsmath, amssymb}
\usepackage{hyperref}
\usepackage{float}
\usepackage{booktabs}
\usepackage{siunitx}
\usepackage{caption}
\usepackage{subcaption}
\usepackage{titlesec} 
\usepackage[utf8]{inputenc}
\usepackage{enumitem}
\raggedbottom
% Bibliography commands
\bibliographystyle{apsrev4-2}
\titleformat{\section}{\bfseries\large\uppercase}{\thesection.}{1em}{}
\titleformat{\subsection}{\bfseries}{\thesubsection}{1em}{}
\begin{document}

\title{QEC strategies against Atom Loss in Neutral Atom Arrays}
\author{Adeniyi Osinowo}
\affiliation{
 Department of Physics and Astronomy,
 University College London,
 Gower Street,
 London,
 WC1E 6BT,
 UK
} 
\date{\today}
\begin{abstract}
% Write abstract based on progress report content
% Quantum error correction (QEC) is a critical component in the development of reliable quantum computing systems.
% Neutral atom arrays have emerged as a promising platform for quantum computing due to their scalability and long coherence times.
% However, atom loss remains a significant challenge that can compromise the effectiveness of QEC strategies.
% This progress report outlines the current state of research on QEC strategies designed to mitigate atom loss in neutral atom arrays.
% Key factors influencing atom loss, relevant error models, and potential QEC strategies are discussed.
% Preliminary results indicate that certain QEC codes show promise in addressing atom loss, but further investigation is needed to optimize their performance.
% Future work will focus on refining error models and implementing QEC strategies to enhance resilience against atom loss.
\end{abstract}
\maketitle

\section{Introduction}% Introduce core concepts and motivation for the review
% - Paragraph: Explain motivation/significance of the topic: 
Since their experimental realization in the early 2000s, neutral atom arrays were initially viewed as scalable but slow and low-fidelity platforms, suitable mainly for quantum simulation rather than quantum computing.
Through rapid progress in recent years, they have achieved significant milestones, including the demonstration of high-fidelity single-qubit gates, two-qubit entangling gates, and the creation of large-scale arrays with hundreds of atoms.
These advancements have made them a leading candidate for scalable quantum computing, offering several advantages over other qubit architectures, including scalability, long coherence times, and the ability to implement high-fidelity parallel entangling gates.
However, atom loss remains a significant challenge that can compromise the effectiveness of QEC strategies on this platform.

% - Paragraph: Define problem statement: Explain how neutral atom arrays work and how atom loss arises, effects of QEC
Neutral atom arrays consist of individual rydberg atoms trapped in optical tweezers or optical lattices, which serve as qubits for quantum computations.
Each atom is manipulated using laser pulses to perform quantum gates, and the state of the qubits can be read out using fluorescence imaging.
During these operations, atoms can be lost due to various factors, including collisions with background gas particles, spontaneous emission, and technical imperfections in trapping mechanisms ~\cite{Bluvstein2024}.
This leads to qubit erasure errors, which differ from standard bit-flip or phase-flip errors in that the location of the error is known, but the state of the qubit is lost.
This type of error can significantly degrade the performance of QEC codes, as it can lead to a loss of logical qubits and reduce the overall fidelity of quantum computations.
% Sentence: Explain why QEC strategies are essential for mitigating the effects of atom loss and ensuring reliable quantum computations in neutral atom arrays 
Effective QEC strategies are essential for maintaining the integrity of quantum information and ensuring reliable quantum computations in the presence of atom loss.

% - Paragraph: Unique challenges of Atom loss for QEC
Standard QEC codes assume a fixed set of physical qubits and errors that flip or phase-flip their states. 
Atom loss removes a qubit entirely, which is fundamentally different from a Pauli error and cannot be corrected by normal syndrome measurements unless the loss is *detected and treated as an erasure*. \cite{PMC}
In common surface codes and related stabilizer QEC codes, each stabilizer involves a set of data qubits. 
If a data qubit is lost, the measurement pattern changes: stabilizers can anticommute with each other, causing flickering syndrome patterns that propagate correlated errors if not properly handled. \cite{Nature}
Once loss events occur, the decoder must incorporate “erasure” information (the locations where qubits have been lost) since this can substantially improve logical error rates if used. 
Decoders that do not use this information perform much worse, while adaptive decoders that incorporate loss can reduce logical failure probability by orders of magnitude. \cite{EmergentMind}
Loss effectively reduces the number of available qubits and can truncate logical code distances. 
If too many qubits are lost before detection, logical qubits may be left underprotected and susceptible to logical errors that QEC cannot correct.

% Paragraph: Explain scope, and approach for reviewing literature and structure of the review
This review focuses on recent research articles, review papers, and experimental studies on QEC strategies addressing qubit losses and erasure errors across various platforms, with the goal of identifying strategies adaptable to neutral atom arrays.
The following sections address the research questions outlined in the introduction.

%   - Identify relevant QEC strategies that can mitigate atom loss in neutral atom arrays
%   - Identify key factors to consider when assessing QEC strategies for atom loss in neutral atom arrays

\section{QEC Background}% Provide background on quantum error correction and its importance for quantum computing
% Paragraph

\section{Modelling atom loss in neutral atom arrays}
% Paragraph: Discuss key factors/challenges of atom loss in neutral atom arrays on QEC performance
% Sentence: Explain how atom loss occurs in neutral atom arrays and the mechanisms behind it (e.g., collisions, spontaneous emission, technical imperfections)
% Sentence: Explain why erasure errors from atom loss are different from standard Pauli errors and how they can impact the performance of QEC codes
% Sentence: Explain the why traditional QEC codes may not be effective against erasure errors and the need for specialized QEC strategies to address this issue
Generally, neutral atom arrays are susceptible to two main sources of errors: depolarizing noise, which leads to standard Pauli errors, and atom loss as well as heralded leakage errors, which lead to erasure errors~\cite{Perrin_2025}.
Traditional QEC codes are primarily designed to correct Pauli errors, and they may not be effective in addressing erasure errors caused by atom loss.
Because the absence of a qubit is easier to detect than standard Pauli errors, erasure errors are easier to correct, but they can still significantly degrade the performance of QEC codes if not properly addressed.
The presence of erasure errors can lead to a significant reduction in the effective code distance, which in turn reduces the error threshold and increases the logical error rate of the QEC code.
Moving atoms to form defect-free arrays or to reconfigure circuits introduces further loss risk because atoms must be accelerated and decelerated within trap arrays, and even slight imperfections or acceleration limits can heat or drop atoms. Algorithms like ATLAS explicitly model this: each move has a *per-move loss probability* that directly reduces the final defect-free fill rate. ([arXiv][2])
State measurement typically involves optical imaging, and the interaction with imaging light can inadvertently kick atoms out of their traps, adding loss during readout. ([Baha Jabarin M.Sc., M.A.][3])

% - Paragraph: Impact on QEC performance
% - Increased logical error rates: Atom loss, if unhandled, turns into correlated syndrome noise or causes faulty stabilizer patterns, raising the logical error rate and degrading the fault-tolerance threshold. 
% Without loss detection and erasure handling, standard QEC approaches may fail even below nominal physical error thresholds. ([Nature][4])
% - Lower fault-tolerance thresholds: The presence of loss imposes additional constraints on the QEC threshold surface. For example, studies show a linear threshold line where allowable depolarizing noise depends on loss probability; high loss can quickly reduce where logical error suppression is feasible. ([Emergent Mind][5])
% - Overhead increases: Managing loss increases overhead through additional ancilla qubits for detection, more complex decoding, and required supply of replacement atoms. These overheads reduce effective code distances and increase circuit complexity per logical operation.
%   - Identify key factors/challenges of atom loss in neutral atom arrays on QEC performance
%   - Identify relevant error models for atom loss in neutral atom arrays
% Paragraph: Discuss relevant error models for atom loss in neutral atom arrays
Studies on atom loss in neutral atom arrays typically employ a combination of atom loss model which accounts for losses during gate operations due to Rydberg excitation and depolarizing noise model, which char-
acterizes imperfections arising from pulse applications and environmental interactions~\cite{Perrin_2025}.
%  - Erasure error models
%  - Stochastic loss models
%  - Loss models based on physical mechanisms of atom loss (e.g., collision-induced loss, spontaneous emission-induced loss)
%  - Hybrid models combining multiple loss mechanisms
%  - Consideration of temporal and spatial correlations in atom loss

\section{Atom loss mitigating QEC strategies}
% Paragraph: Discuss relevant QEC strategies that can mitigate atom loss in neutral atom arrays
%   - Loss Detection Units(LDU) & Teleportation-based LDU
%   - Decoder: Quantum Low-Density Parity-Check (QLDPC) codes with Belief Propagation and OSD Decoders
%   - Decoder: Delayed erasure decoding after loss detection
%   - Dual-rail encoding for loss detection and correction
Previous research has proposed several QEC strategies to mitigate atom loss in neutral atom arrays, including:
\begin{itemize}
    \item Teleportation-based Loss Detection Units (LDU): This approach involves using quantum teleportation to detect and correct for atom loss errors. By teleporting the state of a qubit to a new location, the system can identify if an atom has been lost and take corrective action accordingly~\cite{Perrin_2025}.
    \item Quantum Low-Density Parity-Check (QLDPC) Codes: These codes are designed to handle erasure errors effectively. They can be decoded using belief propagation and ordered statistics decoding (OSD) algorithms, which can improve the performance of QEC in the presence of atom loss~\cite{Roffe_2020}.
    \item Delayed Erasure Decoding: This strategy involves delaying the decoding process until after the loss has been detected, allowing for more accurate error correction based on the known location of the erasure~\cite{Chamberland_2022}.
    \item Dual-rail Encoding: This method encodes quantum information across two physical qubits, allowing for the detection and correction of atom loss by comparing the states of the two qubits~\cite{Chamberland_2022}.
\end{itemize}

% - Paragraph: Mitigation strategies impacting QEC performance
% - Protocols including *leakage-to-erasure conversion* use auxiliary circuits (leakage detection units) to nondestructively sense whether an atom is present. 
% Once a loss is *detected*, the error becomes an *erasure* (known missing qubit), which both reduces decoding complexity and improves thresholds because erasure information can be exploited by decoders. ([APS Journals][6])
% - Improved algorithms like ATLAS plan atom movements while explicitly accounting for per-move loss probabilities to maximize final array fill and reduce initial wasted qubits. 
% Better array layouts directly support QEC by ensuring high occupancy before computation. ([arXiv][2])
% - QEC schemes that modify stabilizer patterns or use combinations of stabilizers (“superchecks”) can maintain valid error-detecting circuits even in the presence of known lost qubits. 
% Decoders that take erasure information as input achieve *lower logical error probabilities* and can even have non-zero thresholds as a function of loss probability. ([Emergent Mind][5])

\section{Conclusion}
% 	- Summary of progress 
This preliminary study has laid the groundwork for a comprehensive analysis of QEC strategies designed to mitigate atom loss in neutral atom arrays. 
Significant progress has been made in identifying key challenges associated with atom loss, relevant error models, and potential QEC strategies. 
Initial implementations of error models and QEC strategies are underway, with promising preliminary results.
Future work will focus on reproducing and extending existing models, as well as exploring additional QEC strategies to enhance the resilience of neutral atom arrays against atom loss.
% 	- Outline and explain difficulties/limitations encountered so far and questions that remain unanswered

\subsection{Assessment methodology}
% Paragraph: Explain how the identified QEC strategies will be assessed for their effectiveness in mitigating atom loss in neutral atom arrays
This study will assess the effectiveness of the identified QEC strategies through a combination of theoretical analysis and numerical simulations.
Theoretical analysis will involve evaluating the error thresholds and logical error rates associated with each QEC strategy under various atom loss scenarios.
Numerical simulations will be conducted to model the performance of each QEC strategy in realistic conditions, taking into account factors such as gate fidelities, coherence times, and the specific error models for atom loss in neutral atom arrays.
The results from both theoretical analysis and numerical simulations will be compared to determine the most effective QEC strategies for mitigating atom loss in neutral atom arrays.

\section*{REFERENCES}
\makeatletter
\let\pre@bibdata\@empty
\makeatother
\bibliography{reference_texts} 

\end{document}